% !TEX program = xelatex
% 厄米多项式与拉盖尔多项式
\documentclass[fontset=windows, chinese, 10pt, xcolor=dvipsnames]{beamer}
% Shared preamble for the chapter 11 standalone lectures.
% Each lecture sets \title and \subtitle after this file is read.

\usepackage{amsmath,amssymb,mathtools}
\usepackage{xeCJK}
\usepackage{zhnumber}
\usepackage{tikz}
\usepackage{tikz-3dplot}
\usetikzlibrary{arrows.meta,calc,positioning}

\usetheme{metropolis}
\usepackage{beamercolorthemeHistNanjing}

\newcommand{\red}[1]{\textcolor{red}{#1}}
\newcommand{\blue}[1]{\textcolor{blue}{#1}}
\newcommand{\green}[1]{\textcolor{green}{#1}}
\newcommand{\gray}[1]{\textcolor{gray}{#1}}
\newcommand{\orange}[1]{\textcolor{orange}{#1}}

\newcommand{\sssub}[1]{ {\mathrm{#1} }}  % script size subscript\\\
\newcommand{\funcd}[2]{ \frac{d{#1}}{d{#2}}} %function full derivative
\newcommand{\funcpd}[2]{ \frac{\partial{#1}}{\partial{#2}}} %function partial derivative
\newcommand{\fnald}[2] { \frac{\delta {#1}}{\delta {#2}}} % functional derivative

\renewcommand{\Re}{\operatorname{Re}}
\renewcommand{\Im}{\operatorname{Im}}
\def\ext{\sssub{ext}}
\def\x{\sssub{x}}
\def\xc{\sssub{xc}}
\def\H{\sssub{H}}
\def\Hx{\sssub{Hx}}
\def\Hxc{\sssub{Hxc}}
\def\c{\sssub{c}}
\def\s{\sssub{s}}

% \def\ext{\sssub{ext}}
% \def\ext{\sssub{ext}}
% \def\ext{\sssub{ext}}
% \def\ext{\sssub{ext}}
% \defsssub{ext}


% definitions of bold faced symbols
\def\brpppp{{\mathbf{r}^{\prime\prime\prime\prime}}}
\def\brppp{{\mathbf{r}^{\prime\prime\prime}}}
\def\brpp{{\mathbf{r}^{\prime\prime}}}
\def\brp{{\mathbf{r}^{\prime}}}
\def\bzp{{\mathbf{z}^{\prime}}}
\def\bxp{{\mathbf{x}^{\prime}}}
\def\tp{{{t}^{\prime}}}
\def\tpp{{{t}^{\prime\prime}}}
\def\tppp{{{t}^{\prime\prime\prime}}}

\def\tbr{{\tilde{\mathbf{r}}}}
\def\bk{{\mathbf{k}}}
\def\br{{\mathbf{r}}}
\def\bp{{\mathbf{p}}}
\def\bq{{\mathbf{q}}}
\def\bv{{\mathbf{v}}}
\def\bz{{\mathbf{z}}}
\def\mbf{{\mathbf{f}}}

\def\bx{{\mathbf{x}}}
\def\bR{{\mathbf{R}}}
\def\bM{{\mathbf{M}}}
\def\bP{{\mathbf{P}}}
\def\bT{{\mathbf{T}}}
\def\bK{{\mathbf{K}}}
\def\bA{{\mathbf{A}}}
\def\bB{{\mathbf{B}}}
\def\bD{{\mathbf{D}}}
\def\bE{{\mathbf{E}}}
\def\bF{{\mathbf{F}}}
\def\bG{{\mathbf{G}}}
\def\bH{{\mathbf{H}}}
\def\bI{{\mathbf{I}}}
\def\bJ{{\mathbf{J}}}
\def\bP{{\mathbf{P}}}
\def\bS{{\mathbf{S}}}
\def\bV{{\mathbf{V}}}
\def\bX{{\mathbf{X}}}
\def\bY{{\mathbf{Y}}}
\def\dd{{\mathop{}\!\mathrm{d}}}
\def\dr{{\!\!\dd\br\,}}
\def\dx{{\!\!\dd\bx\,}}
\def\dy{{\!\!\dd\by\,}}
\def\dk{{\!\!\dd\bk\,}}
\def\dG{{\!\!\dd\bG\,}}


\def\half{{\frac{1}{2}}}
\def\tr{{\mathrm{tr}}}


\def\Arg{{\mathrm{Arg}\;}}
\def\arg{{\mathrm{arg}\;}}
\def\Res{{\mathrm{Res}\;}}

% unit-vector basis (bold under XeLaTeX; \mathbf can look non-bold)
\newcommand{\ue}{\boldsymbol{e}}
\newcommand{\un}{\hat{\boldsymbol{n}}}


\author{罗凯}
\institute{南京理工大学物理学院}
\date{\zhdigits{\the\year},\zhganzhinian{\the\year}年秋}

\usebackgroundtemplate{%
\begin{tikzpicture}[remember picture,overlay]
  \node[opacity=0.05, at=(current page.center)] {
    \includegraphics[width=0.6\paperwidth]{njust-logo.jpeg}
  };
\end{tikzpicture}}


\title{特殊函数简介}
\subtitle{第四讲：厄米多项式与拉盖尔多项式}

\begin{document}

\begin{frame}
  \titlepage
\end{frame}

\begin{frame}{本讲提纲}
  勒让德多项式来自球面上的角向方程。谐振子和氢原子再产生两类多项式，它们在无穷区间上带权重正交。
  \begin{enumerate}
    \item 谐振子要求厄米多项式在有限项截断，能级是 $\hbar\omega(n+1/2)$；
    \item 罗德里格斯公式、生成函数和权重 $e^{-\xi^2}$；
    \item 氢原子的径向函数是广义拉盖尔多项式，能级只依赖 $n$。
  \end{enumerate}
\end{frame}

\section{厄米多项式}

\begin{frame}{谐振子}
  \[
    -\frac{\hbar^2}{2m}\frac{\mathrm{d}^2\psi}{\mathrm{d}x^2}
    +\frac12 m\omega^2 x^2\,\psi
    =E\psi.
  \]
  令 $\xi=x\sqrt{m\omega/\hbar}$，$\varepsilon=2E/(\hbar\omega)$，
  \[
    -\frac{\mathrm{d}^2\psi}{\mathrm{d}\xi^2}+(\xi^2-\varepsilon)\psi=0.
  \]
  $\xi\to\infty$ 时可积的解按 $e^{-\xi^2/2}$ 衰减。令 $\psi=v(\xi)\,e^{-\xi^2/2}$，
  \[
    v''-2\xi v'+(\varepsilon-1)v=0.
  \]
\end{frame}

\begin{frame}{级数必须截断}
  在 $\xi=0$ 把 $v$ 写成幂级数。递推每一步把次数提高 $2$，一般项像 $\xi^{2k}/k!$ 再乘一个不趋于零的因子。这样 $v$ 比 $e^{\xi^2/2}$ 长得还快，$\psi$ 在无穷远处不可积。
  \vspace{0.4em}
  级数在有限项截断的条件是 $\varepsilon-1=2n$，$n=0,1,2,\ldots$。能级因此是
  \begin{equation*}
    \boxed{\displaystyle
      E_n=\hbar\omega\Bigl(n+\frac12\Bigr).}
  \end{equation*}
\end{frame}

\begin{frame}{罗德里格斯公式}
  \begin{equation*}
    \boxed{\displaystyle
      H_n(\xi)
      =(-1)^n e^{\xi^2}
        \frac{\mathrm{d}^n}{\mathrm{d}\xi^n}e^{-\xi^2}.}
  \end{equation*}
  $H_0=1$，$H_1=2\xi$。二阶导数
  \[
    \frac{\mathrm{d}^2}{\mathrm{d}\xi^2}e^{-\xi^2}=(4\xi^2-2)e^{-\xi^2},
  \]
  所以 $H_2=4\xi^2-2$。再求一次，$H_3=8\xi^3-12\xi$。
\end{frame}

\begin{frame}{核对 $H_2$}
  $H_2'=8\xi$，$H_2''=8$。厄米方程是 $H_n''-2\xi H_n'+2n H_n=0$。$n=2$ 时
  \begin{align*}
    &8-2\xi\cdot 8\xi+4(4\xi^2-2)\\
    &\qquad=8-16\xi^2+16\xi^2-8=0.
  \end{align*}
  这就是 $\varepsilon-1=2n$ 时 $v$ 的方程。
\end{frame}

\begin{frame}{生成函数}
  \[
    e^{2\xi t-t^2}
    =\sum_{n=0}^{\infty} H_n(\xi)\frac{t^n}{n!}.
  \]
  对 $t$ 求导，比较系数，
  \[
    H_{n+1}(\xi)=2\xi H_n(\xi)-2n H_{n-1}(\xi).
  \]
  $H_0$ 和 $H_1$ 之后，后面的多项式逐个写出来。
\end{frame}

\begin{frame}{正交与波函数}
  权重是 $e^{-\xi^2}$，
  \[
    \int_{-\infty}^{\infty}
      e^{-\xi^2} H_m(\xi) H_n(\xi)\,\mathrm{d}\xi
    =\sqrt{\pi}\,2^n n!\,\delta_{mn}.
  \]
  归一化的束缚态是
  \[
    \psi_n(x)
    =(2^n n!)^{-1/2}
      \Bigl(\frac{m\omega}{\pi\hbar}\Bigr)^{1/4}
      H_n(\xi)\,e^{-\xi^2/2},
  \]
  $\xi=x\sqrt{m\omega/\hbar}$。基态 $n=0$ 没有节点，$H_0=1$，概率密度是一个高斯。
\end{frame}

\section{拉盖尔多项式}

\begin{frame}{氢原子的径向方程}
  波函数写成 $R(r)Y_{\ell m}$。令 $u(r)=r R(r)$，
  \[
    -\frac{\hbar^2}{2m}u''
    +\Biggl(-\frac{e^2}{r}
      +\frac{\hbar^2\ell(\ell+1)}{2m r^2}\Biggr)u
    =E u.
  \]
  束缚态 $E<0$。令 $\kappa=\sqrt{-2mE}/\hbar$，$\rho=2\kappa r$，以及
  \[
    n=\frac{m e^2}{\hbar^2\kappa}.
  \]
  这里的 $n$ 先只是一个参数。
\end{frame}

\begin{frame}{能级}
  $\rho\to\infty$ 时 $u\sim e^{-\rho/2}$，原点附近 $u\sim\rho^{\ell+1}$。提出这两个因子，
  \[
    u(r)=\rho^{\ell+1}e^{-\rho/2} w(\rho),
  \]
  $w$ 满足广义拉盖尔方程。可积性要求级数截断：$n$ 是正整数，并且 $n\geq\ell+1$。能级是
  \begin{equation*}
    \boxed{\displaystyle
      E_n=-\frac{m e^4}{2\hbar^2}\frac{1}{n^2}.}
  \end{equation*}
  高斯单位制下大约是 $-13.6\,\mathrm{eV}/n^2$。能量只依赖 $n$。
\end{frame}

\begin{frame}{简并度}
  固定 $n$ 时，$\ell=0,1,\ldots,n-1$，每个 $\ell$ 又有 $2\ell+1$ 个 $m$。
  \[
    \sum_{\ell=0}^{n-1}(2\ell+1)=n^2.
  \]
  这 $n^2$ 重简并来自库仑势 $1/r$。微扰一般会把这 $n^2$ 个态拆开。
\end{frame}

\begin{frame}{广义拉盖尔多项式}
  \begin{equation*}
    L_n^{(\alpha)}(x)
    =\frac{e^{x} x^{-\alpha}}{n!}
      \frac{\mathrm{d}^n}{\mathrm{d}x^n}
      \bigl(e^{-x} x^{n+\alpha}\bigr),
    \qquad\alpha>-1.
  \end{equation*}
  $\alpha=0$ 时是普通拉盖尔多项式 $L_n$。前几个是
  \[
    L_0^{(\alpha)}=1,
    \qquad
    L_1^{(\alpha)}=1+\alpha-x.
  \]
  权重 $e^{-x}x^{\alpha}$ 下，
  \[
    \int_0^{\infty}
      e^{-x}x^{\alpha}
      L_n^{(\alpha)} L_m^{(\alpha)}\,\mathrm{d}x
    =\frac{\Gamma(n+\alpha+1)}{n!}\,\delta_{mn}.
  \]
\end{frame}

\begin{frame}{径向波函数}
  \[
    w(\rho)=L_{n-\ell-1}^{(2\ell+1)}(\rho),
    \qquad
    \rho=\frac{2r}{n a_0},
    \qquad
    a_0=\frac{\hbar^2}{m e^2}.
  \]
  $n-\ell-1$ 是径向节点数。
  \vspace{0.35em}
  基态 $n=1$、$\ell=0$ 没有节点，$L_0^{(1)}=1$。径向概率密度 $r^2|R|^2$ 在 $r=a_0$ 取最大。
\end{frame}

\begin{frame}{这一讲留下的公式}
  \[
    H_n(\xi)
    =(-1)^n e^{\xi^2}\frac{\mathrm{d}^n}{\mathrm{d}\xi^n}e^{-\xi^2},
    \qquad
    E_n=\hbar\omega\Bigl(n+\frac12\Bigr).
  \]
  \[
    E_n=-\frac{m e^4}{2\hbar^2 n^2},
    \qquad
    w(\rho)=L_{n-\ell-1}^{(2\ell+1)}(\rho).
  \]
  厄米多项式的权重是 $e^{-\xi^2}$，拉盖尔多项式的权重是 $e^{-x}x^{\alpha}$。氢原子的 $n$ 能级简并度是 $n^2$。
\end{frame}

\end{document}
